\documentclass[main.tex]{subfiles}
%subfile: ch4.tex
\begin{document}

	
\chapter{Works of H. J. K. Junnila} \label{ch4}

\section{Submetacompact and Subparacompact Spaces} \label{sec:4.1}

{\definition{} \label{def:4.1.1}
Let $\mathscr{U}=\{U_\alpha: \alpha\in\Delta\}$ and $\mathscr{V}$ be covers of $X$. 
We extend \uwave{point-star refinement} in Defition \ref{def:1.1.6}. 

\textup{(i)} $\mathscr{V}$ is a 
\uwave{point-star}
\index{point-star refinement| \S\ref{def:4.1.1}\hfill}
$($\uwave{local-star}\index{local-star refinement| \S\ref{def:4.1.1}\hfill}$)$ 
\uwave{refinement} of $\mathscr{U}$ 
at $x \in X$ if there exists $\alpha\in\Delta$ $($and $G \in\operatorname{top}(x))$
such that $\operatorname{st}(x, \mathscr{V}) \subset U_\alpha$
$(\operatorname{st}(G, \mathscr{V}) \subset U_\alpha)$. 

\textup{(ii)} $\mathscr{V}$ is a 
\uwave{point-star $\dot{F}$-refinement}
\index{point-star $\dot{F}$-refinement| \S\ref{def:4.1.1}\hfill}
of $\mathscr{U}$ at $x$ if there exists 
$F \in [\Delta]^{<\omega}$ such that $x \in \bigcap_{\alpha \in F} U_\alpha$ 
and $\operatorname{st}(x, \mathscr{V}) \subset \bigcup_{\alpha\in F}U_\alpha$.

\textup{(iii)} $\mathscr{V}$ is a 
\uwave{point-star} $($\uwave{local-star}, 
\uwave{point-star $\dot{F}$}$)$ \uwave{refinement}  of $\mathscr{U}$ if $\mathscr{V}$ is a point-star $($local-star, point-star $\dot{F})$ refinement
of $\mathscr{U}$ at all $x \in X$.

\textup{(iv)} A sequence $\{\mathscr{V}_n: n \in \omega\}$ of covers of $X$ is a 
\uwave{point-star} $($\uwave{local-star}, 
\uwave{point-star $\dot{F}$}$)$ \uwave{refining sequence} 
\index{point-star refining sequence| \S\ref{def:4.1.1}\hfill}
\index{local-star refining sequence| \S\ref{def:4.1.1}\hfill}
\index{point-star $\dot{F}$ refining sequence| \S\ref{def:4.1.1}\hfill}
for $\mathscr{U}$ if for 
each $x \in X$ there exists $n < \omega$ such that $\mathscr{V}_n$ is a point-star $($local-star, point-star $\dot{F})$ refinement of $\mathscr{U}$ at $x$.
}

{\theorem{} \label{th:4.1.2} 
\textup{([Junnila 1978]\textsuperscript{\cite{Junnila1978}})}
$X$ is $\lambda$-submetacompact if and only if for every open $\lambda$-cover $\mathscr{U}$ of $X$, 
there exists a point-star $\dot{F}$-refining sequence for $\mathscr{U}$ 
by semi-open covers of $X$.
}

\begin{proof}
$(\Rightarrow)$. If $\{\mathscr{V}_n: n < \omega\}$ is a $\theta$-sequence of open refinement of a 
cover $\mathscr{U}$, then $\{\mathscr{V}_n: n < \omega\}$ is a point-star $\dot{F}$-refining
sequence for $\mathscr{U}$ by semi-open covers of $X$.

$(\Leftarrow)$. Suppose that every open $\lambda$-cover has a point-star $\dot{F}$-refining 
sequence by semi-open covers, and that $\mathscr{U} = \{U_\alpha: \alpha < \lambda\}$ 
be an open cover of $X$. We use induction on $n < \omega$ to define open covers 
$\mathscr{V}(t) = \{V(t, \beta): \beta < \lambda \cdot 2\}$ of $X$ for 
$t \in \omega^n \subset \omega^{<\omega}$ such that:
 
\begin{equation}
V(t, \alpha) \subset U_\alpha \textup{ and } V(t, \lambda+\alpha) \subset U_\alpha 
\textup{ for each } \alpha < \lambda.      \label{eq:4.2.2.1}
\end{equation}

\noindent Let $V(\emptyset, \alpha) = \emptyset$ and $V(\emptyset, \lambda + \alpha) = U_\alpha$ 
for each $\alpha < \lambda$. 
Let $\mathscr{V}(\emptyset) = \{V(\emptyset, \beta): \beta < \lambda \cdot 2\}$. 
Suppose that $n < \omega$ and that we have defined the open covers 
$\mathscr{V}(s) = \{V(s, \beta): \beta < \lambda \cdot 2\}$ for each $s \in \omega^{n-1}$.
For each $s \in \omega^{n-1}$, $\mathscr{V}(s)$ has a point-star $\dot{F}$-refining sequence $\{\mathscr{L}_{s \oplus k}\}_{k < \omega}$ by semi-open covers of $X$. 
We now define, for all $s \in \omega^{n-1}$, $k < \omega$, and $\alpha < \lambda$:

\begin{gather}
V(s \oplus k, \alpha) = V(s, \alpha) \bigcup \left( U_\alpha \bigcap 
\text{st}\left( X \setminus \bigcup_{\beta \ne\lambda+ \alpha} V(s, \beta), 
\mathscr{L}_{s \oplus k} \right)^\circ \right),      \label{eq:4.2.2.2} \\
V(s \oplus k, \lambda + \alpha) = U_\alpha \bigcap 
\left( \bigcup_{\beta > \lambda + \alpha} V(s, \beta) \right) 
\bigcap \text{st}\left( X \setminus \bigcup_{\beta < \lambda + \alpha} V(s, \beta), 
\mathscr{L}_{s \oplus k} \right)^\circ.   \label{eq:4.2.2.3}
\end{gather}
Since every element $t \in \omega^n$ can be uniquely represented in the form $s \oplus k$ 
for some $s\in \omega^{n-1}$ and $k < \omega$, 
the above formulas define the families $\mathscr{V}(t) = \{V(t, \beta): \beta < \lambda \cdot 2\}$ 
for $t \in \omega^n$. To complete the induction, it remains to verify that the families 
$\mathscr{V}(t)$, $t \in \omega^n$, are open covers of $X$.
It is easy to see that these families consist of open subsets. 
To show that they cover $X$, let $t = s \oplus k \in \omega^n$ and let $x \in X$.
Since $\mathscr{V}(s)$ covers $X$, let 
$\beta_0 = \min \{ \alpha < \lambda \cdot 2: x \in \mathscr{V}(s, \beta) \}.$
We have
\begin{equation}
x \in V(s, \beta_0)  \textup{ and }
x \notin \bigcup_{\beta < \beta_{0}} V(s, \beta).   \label{eq:4.2.2.4}
\end{equation}
If $\beta_0 < \lambda$, then $V(s, \beta_0)\subset V(t, \beta_0)$. 
So $x \in V(t, \beta_0)$.
Assume that $\beta_{0} \ge \lambda$, then $\beta_{0} = \lambda + \alpha_{0}$ 
for some $\alpha_0 < \lambda$. Then $x \in V(s, \lambda + \beta_0) \subset U_\alpha.$
\begin{enumerate}
\item[(a)] If $x \notin \bigcup_{\beta > \lambda + \alpha_{0}} V(s, \beta)$, 
	then by \eqref{eq:4.2.2.4} and \eqref{eq:4.2.2.2}
	\[ x \in U_{\alpha_{0}} \cap \operatorname{st}\left(x \setminus 
	\bigcup_{\beta \ne \lambda + \alpha_{0}} V(s, \beta), 
	\mathscr{L}_{s \oplus k}\right)^{\circ} 
	\subset V(t, \alpha_{0}) \]
\item[(b)] If $x \in \bigcup_{\beta > \lambda + \alpha_{0}} V(s, \beta)$, 
	then by \eqref{eq:4.2.2.4} and \eqref{eq:4.2.2.3}
	\[ x \in U_{\alpha_{0}} \cap \left(\bigcup_{\beta > \lambda + \alpha_{0}} U(s, \beta)\right) 
	\cap \operatorname{st}\left(X- \bigcup_{\beta < \lambda + \alpha_{0}} V(s, \beta), 	\mathscr{L}_{s \oplus k}\right)^{\circ} 
	\subset V(t, \alpha_{0}) \]
\end{enumerate}
Thus $\mathscr{V}(t)$ covers $X$. 
By  \eqref{eq:4.2.2.1}, we have now defined the open refinements 
$\mathscr{V}(t)$ for $t \in \omega^{<\omega}$.
Since $\omega^{<\omega}$ is countable, tocomplete the proof, it suffices to show that
\begin{equation}
	\{\mathscr{V}(t): t \in \omega^{<\omega}\} \text{ is a } \theta\text{-sequence.}  \label{eq:4.2.2.5}
\end{equation}
\noindent Let $x \in X$. There is a sequence $\{s_n\}$ such that 
$s_0 = \emptyset\in \omega^0, s_1 = s_0\oplus k, 
\cdots, s_{n+1} = s_n \oplus k$ where $s_n \in \omega^n, k < \omega$, 
and such that $\mathscr{L}_{s_{n+1}}$ is a point-star $\dot{F}$-refinement of
$\mathscr{V}(s_n)$ at $x$ for each $n < \omega$.
For each $n < \omega$, let $A_n$ be a finite subset of $\lambda\cdot 2$ such that 
$x \in \bigcap_{\beta \in A_n} V(s_n, \beta)$ and
\begin{equation}
	\operatorname{st}(x, \mathscr{L}_{s_{n+1}}) \subset \bigcup_{\beta \in A_n} V(s_n, \beta),
\forall \, n < \omega.   \label{eq:4.2.2.6}
\end{equation}
\noindent Let $\beta_n = \max A_n$ and $\beta_m = \min\{\beta_n : n < \omega\}$. We have
\begin{equation}
\text{If } \beta \ge \lambda, \text{ then } x \notin V(s_{m+1}, \beta). \label{eq:4.2.2.7}
\end{equation}

\noindent We proceed with two parts.

(a) If $\beta \ge \lambda$ and $\beta > \beta_m$, then $x \notin V(s_{m+1}, \beta)$.

\noindent The assumption of (a) implies that $\gamma \le \beta_m < \beta$ for all $ \gamma\in A_m$. 
Then by \eqref{eq:4.2.2.6}, 
$\operatorname{st}(x, \mathscr{L}_{s_{m+1}}) \subset \bigcup_{\gamma < \beta} V(s_m, \gamma)$. 
Thus: $x \notin \operatorname{st}
\left(X - \bigcup_{\gamma < \beta} V(s_m, \gamma), \mathscr{L}_{s_{m+1}}\right)$.
It follows from \eqref{eq:4.2.2.3} that $x\notin V(s_{m+1}, \beta)$, (a) is true.

(b) $\beta_m < \lambda$.

\noindent For if $\beta_m \ge \lambda$, 
$\beta > \beta_m \ge \lambda$ for all $\beta > \beta_{m+2}$
since $\beta_{m+2} \ge \beta_m$.
So by (a), $x\notin V(s_{m+1}, \beta)$. 
Hence $x\notin \bigcup_{\beta> \beta+2}V(s_{m+1}, \beta)$.
Since $\beta_{m+2} \ge \lambda$, by \eqref{eq:4.2.2.3},
$x\notin V(s_{m+2}, \beta_{m+2})$. But $\beta_{m+2} \in A_{m+2}$,
this contradiction, showsthat part (b) is true.
Then \eqref{eq:4.2.2.7} follows from (a) and (b). We say that

\begin{equation}
	\mathscr{V}(s_{m+1}) \text{ is point-finite at } x. \label{eq:4.2.2.8}
\end{equation}
\noindent If this is true, then \eqref{eq:4.2.2.5} is true and we complete the proof. 

Now, we prove \eqref{eq:4.2.2.8}.
Let $C= \operatorname{ind}(\mathscr{V}(s_{m+1}))_x$. By \eqref{eq:4.2.2.7}, $C = \{\alpha < \lambda : x \in V(s_{m+1}, \alpha)\}$.For $n\le m$, let $B_n = \{\alpha < \lambda: \lambda + \alpha\in A_n\}$. Since each
$B_n $ is finite, it suffices to show that $C \subset \bigcup_{n=1}^{m} B_n$. 
Let $\alpha \in C$. Since $x\notin \emptyset = V(s_0, \alpha)$, there exists $p$ such that 
$0 \le p \le m$ and $x \in \mathscr{V}({s_{p+1}}, \alpha)\setminus  \mathscr{V}{s_{p}}, \alpha)$.
By \eqref{eq:4.2.2.2}, 
$x \in \operatorname{st}(X \setminus \bigcup_{\beta \ne \lambda + \alpha}\mathscr{V}(s_p, \beta),
	 \mathscr{L}_{s_{p+1}})$.
Let $y \in \operatorname{st}(x, \mathscr{L}_{s_{p+1}})\subset 
\bigcup_{\beta\ne \lambda + \alpha}\mathscr{V}(s_p, \beta)$
be such that $x \in \operatorname{st}(y, \mathscr{L}_{s_{p+1}})$. Then
 $y \in \operatorname{st}(x, \mathscr{L}_{s_{p+1}})\subset\bigcup_{\beta\in A_p}\mathscr{V}(s_p, \beta)$.
 Suppose $y \in V(s_p, \gamma)$ for some $\gamma\in A_p$.
 Then $\lambda + \alpha= \gamma \ in A_p$, so $\alpha \in B_p$.
We have show that $C \subset \bigcup_{n=1}^{m} B_n$, and hence 
$\mathscr{V}(s_{m+1})$ is point-finite at $x$.
\end{proof}

{\corollary{} \label{cor:4.1.3} 
For any space $X$, the following are equivalent:

\textup{(i)} $X$ is $\lambda$-submetacompact.

\textup{(ii)}  Eveery open $\lambda$-cover of $X$ has a $\theta$-sequence of  open reinements.

\textup{(iii) ([Worrell and Wicke 1965]\textsuperscript{\cite{WorrellWicke}})} Every open $\lambda$-cover of $X$ has 
a pointwise $W$-refining sequence by open covers. \hfill\qedsymbol
}

{\lemma{} \label{lem:4.1.4}
If every interior-preseving open $\lambda$-cover of $X$ has a point-star $\dot{F}$  
refining sequence of interior-preseving open covers, 
then every interior-preseving open $\lambda$-cover has a $\theta$-sequence of open refinements.
}


\begin{proof}
In the proof of theorem \ref{th:5.1.2}, if the open covers $\mathscr{U}$ and each $\mathscr{U}_{t}, t \in \omega^{<\omega}$, are open interior-preseving, then it is easy to see that the covers 
$\mathscr{V}_{(t)}, t \in \omega^{<\omega}$, defined by (2) and (3) of the proof 
will also be interior-preseving open.
\end{proof}

{\theorem{} \label{th:4.1.5}
$X$ is $\lambda$-submetacompact iff every increasing open $\lambda$-cover has a 
$\theta$-sequence of open refinements.
}

\begin{proof} 
($\Leftarrow$). Clear.

($\Rightarrow$). Assume that every increasing open $\lambda$-cover has a 
$\theta$-sequence of open refinements. 
Let $\lambda \ge \omega$. For each $\kappa < \lambda$, denote by $\mathcal{P}_k$ 
the proposition: "every open $\kappa$-cover of $X$ has a pointwise $W$-refining 
sequence by open covers of $X$". By Corollary \ref{th:5.1.3}, $X$ is $\kappa$-submetacompact iff $\mathcal{P}_\kappa$. We use transfinite induction to show that $\mathcal{P}_\kappa$ holds for 
each $\kappa \le \lambda$. If $\kappa$ is finite, $\mathcal{P}_\kappa$ is trivially true.
Let $\kappa \le \lambda$ such that $\mathcal{P}_\mu$ holds for each $\mu < \kappa$.
To show $\mathcal{P}_\kappa$ holds, let $\mathscr{U} = \{U_\alpha : \alpha < \kappa\}$ be an open 
cover of $X$. The family $\mathscr{U}^*=\{\bigcup_{\gamma\le\alpha}U_\gamma: \alpha < \kappa\}$ 
is an increasing open $\kappa$-cover of $X$. 
It has a $\theta$-sequence $\{\mathscr{V}_n: n < \omega\}$ of open refinements. 
We may assume that each $\mathscr{V}_n = \{V(n, \alpha) : \alpha < \kappa\}$ 
such that $V(n, \alpha) \subset \bigcup_{\gamma \le \alpha} U_\gamma$ for all $\alpha < \kappa$. 
For $n < \omega, \alpha < \kappa$, let
$$\mathscr{U}(n, \alpha) = \{U_\gamma : \gamma \le \alpha\} \cup 
\{\bigcup_{\gamma > \alpha} V(n, \gamma)\}.$$
\noindent Then each $\mathscr{U}(n, \alpha)$ is an open cover of $X$ with power less than $\kappa$.
By the induction hypothesis, it has a pointwise W-refining sequence 
$\{\mathscr{W}_l(n, \alpha): l < \omega\}$ by open covers of $X$.
For $n, h < \omega$, let $F_{n,0} = \emptyset$ and 
$$F_{n,h} = \{x \in X: |\operatorname{ind}(\mathscr{V}_n)_x| \le h\} \text{ for } h \ge 1.$$
Then each $F_{n,h}$ is closed.
 For each $y \in F_{n,h}$, let $\delta(y, n) = \max(\text{ind}(\mathscr{V}_n)_y)$.
We have 

(1) $y \in V(n, \delta(y, n))$ and $y\notin \bigcup_{\gamma > \delta(y, n)} V(n, \gamma).$

\noindent For each $l = 0, 1, \dots, h$, let $W_l(n, y)$ be a member of $\mathscr{W}_l(n, \delta(y, n))$ containing $y$.
Let 
$$H(n, h, y) = V(n, \delta(y, n)) \cap \bigcap_{l=0}^{h} W_l(n, y).$$ 
\noindent Then $H(n, h, y)$ are open sets containing $y$.
Let 
$$\mathscr{H}_{n,h} = \{H(n, h, y) : y \in F_{n,h}\} \cup \{X \setminus F_{n,h}\}.$$
\noindent Then $\mathscr{U}_{n,h}$ are open covers of $X$.
We have

(2) $\{\mathscr{H}_{n,h}\}_{n,h < \omega}$ is a pointwise W-refining sequence for $\mathscr{U}$, so $\mathcal{P}_\kappa$ holds.

\textsc{Proof of (2)}. Let $x \in X$. There exists $m$ such that $\text{ind}(\mathscr{V}_m)_x = \{\alpha_1, \alpha_2, \dots, \alpha_p\}$ is finite. Then $x \in F_{m,p}$. For each $i = 1, \dots, p$, since 
$\{\mathscr{W}_l(m, \alpha_i): l < \omega\}$ is a pointwise W-refining sequence for 
$\mathscr{U}(m, \alpha_i)$, 
there exists $l_i < \omega$ and a finite set $D_i \subset \alpha_i$ such that:

(3) $(W_{l_i}'(m, \alpha_i))_x$ refines $\{U_\gamma : \gamma \in D_i\} \cup
\{\bigcup_{\gamma > \alpha_i}V(m, \gamma)\}$.

\noindent Let $\mathscr{U}' = \{U_\gamma : \gamma \in \bigcup_{i=1}^{p} D_i\}$ and 
$r = \max\{p,  l_1, \dots, l_p\}$.

(4) $(\mathscr{H}_{m,r})_x$ refines $\mathscr{U}'$, so (2) holds.

\textsc{Proof.} Let $H \in (\mathscr{H}_{m,r})_x$, then $x\in H\cap F_{m,r}$.
Hence there exists $y \in F_{m,r}$ such that $H= H(m, r, y)$.
Since $x\in V(m, \delta(y, m))$, there exists $j \le p$ such that $\delta(y, m) = \alpha_j$.
Also, $x\in\bigcap_{l=0}^r W_l(m, y)\subset W_{l_j}(m, y)$. It follows from (3) and (1), there exists
$\gamma \in D_j$ such that $W_{l_j}(m, y) \subset U_\alpha$,
$H \subset U_{\alpha} \in \mathscr{U}'$, thus (4) holds. This completes the proof that $\mathcal{P}_\kappa$ holds.
\end{proof}

{\lemma{} \label{lem:4.1.6}
Let $\mathscr{U}$ be an open cover of $X$.

\textup{(i)} If $\mathscr{U}$ has an interior-preseving open point (local)-star refinement, 
then $\mathscr{U}$ has a closure-preserving closed refinement 
$\mathscr{F}$ (with $\bigcup \mathscr{F}^\circ = X$).

\textup{(ii)} If $\mathscr{U}$ has a point (local) star refining sequence of interior-preseving 
open covers of $X$, then $\mathscr{U}$ has a $\sigma$-closure-preserving closed 
refinement $\mathscr{F}$ (with $\bigcup \mathscr{F}^\circ = X$).
}

* **Proof of Lemma 4.1.6 (i):** Suppose $\mathscr{U} = \{V_\alpha : \alpha \in \Gamma\}$ has a point (local) star refining sequence of interior-preseving open covers of $X$. For each $n < \omega, \alpha \in \Gamma$, let $F_{n,\alpha} = \{x \in X : \operatorname{st}(x, \mathscr{U}_n) \subset V_\alpha\}$. Then $F_{n,\alpha}$ is closed for all $\alpha \in \Gamma$.


* To see that each $\mathscr{F}_n = \{F_{n,\alpha} : \alpha \in \Gamma\}$ is ep closed, let elements satisfy the required intersection and containment properties.


* $\mathscr{V}$ is a $\sigma$-ep closed refinement of $\mathscr{U}$ since $\bigcup_{n<\omega} \mathscr{V}_n = X$, and for any $x \in X$, there is $m$ and $G \in \mathscr{U}_m$ such that $\operatorname{st}(x, \mathscr{U}_m) \subset V_\alpha$ for some $\alpha$, ensuring $x \in F_{m,\alpha}$ and $\bigcup \mathscr{V} = X$.

{\theorem{} \label{th:4.1.7}
For any space $X$, the following statements are equivalent:

\textup{(i)}  $X$ is $\sigma$-submetacompact.

\textup{(ii)}  Every directed open $\sigma$-cover of $X$ has a $\sigma$-cp closed refinement.

\textup{(iii)}  Every interior-preseving directed open $\sigma$-cover of $X$ has a $\sigma$-cp closed refinement.

\textup{(iv)} Every interior-preseving open $\sigma$-cover of $X$ has a pointwise $W$-refining sequence of interior-preseving open covers of $X$.

\textup{v)}  Every interior-preseving open $\sigma$-cover of $X$ has a point-star $\sigma$-refining sequence of interior-preseving open covers of $X$.

\textup{(vi)}  Every interior-preseving directed open $\sigma$-cover of $X$ has a point-star refining sequence of interior-preseving open covers of $X$.
}

\begin{proof} 
 Let $\mathscr{U}$ be a directed open $\sigma$-cover of $X$.


* By (i) and Theorem 3.1.5, there is a countable closed cover $\{F_n\}$ of $X$ such that each subspace has a point-finite open refinement $\mathscr{V}_n$ in $F_n$.


* Each $\mathscr{V}_n$ is an interior-preseving open point-star refinement in $F_n$.


* By Lemma 4.1.6, each has a cp closed refinement in the closed subspace $F_n$, providing a $\sigma$-cp closed refinement of $\mathscr{U}$.




* **Additional Proof Connections:**
* Relationships and implications between (iii), (iv), and subsequent conditions are established using prior lemmas and definitions.
\end{proof}{proof} 

\subfile{ch4_2.tex}
\end{document}
